\Section{Introduction}
\label{sec:intro}

The Move Prover (\MVP)~\cite{DillGPQXZ22} is a formal verification tool for Move smart contracts \cite{blackshear2020resources, MoveOnAptosBook}, designed for routine use during code development, with running times comparable to type checkers and linters.
\MVP is deployed at Aptos~\cite{ParkZGXGCLC24, AptosFrameworkBook} for formal verification of core protocol logic such as staking, metering, code deployment, and auxiliary data structures that underpin the framework.

%  Recently, \MVP was extended to support Move's higher-order functions and dynamic dispatch ~\cite{grieskamp2025movehigherorder}, including behavioral predicates over function values and state labels naming intermediate memory states~\cite{GrieskampEtAl26HO}.
%
A practical bottleneck in program verification is the cost of \emph{writing} the specifications that drive verification.
The properties developers want to establish are typically high-level: that a call chain aborts only under stated preconditions, that a global storage invariant is preserved, that funds are conserved across a transfer.
While \MVP can verify unspecified callees by inlining their bodies, this does not scale: modular verification requires explicit per-function pre- and postcondition specifications.
This detail is mechanical and tedious to write, dominates the cost of using a verifier, and is a significant barrier to wider adoption.

\emph{Specification inference} is a promising approach. Inference has a long history \cite{DaikonTSE01} and has gained steam again with LLMs \cite{PeiEtAl23, KamathEtAl23}.
Coding agents like Claude Code \cite{ClaudeCode} open further avenues: they expose tools to the AI (e.g. via the Model Context Protocol, MCP) and provide a workflow in which AI and developers collaborate on code changes. And models are significantly more powerful than three years ago.

We present early work on a novel combination of (a) mechanical spec inference via weakest-precondition analysis (\WP)~\cite{dijkstra1975wp,BL05} and (b) agentic inference guided by a set of skills (English instructions), integrated as a plugin into Claude Code \cite{ClaudeCodePlugins}.
Our goal is a methodology in which AI, mechanical deduction, and users collaborate on writing and verifying Move specifications.
This paper defines some of the building blocks toward that goal, focusing on spec inference.

By reusing \MVP's reference-elimination pass~\cite{DillGPQXZ22} (also adopted by Aeneas~\cite{Aeneas}), \WP propagates over a reference-free intermediate form on which classical Dijkstra-style reasoning applies directly, and recovers abort conditions, modifies clauses, and consequence-style postconditions from the bytecode.

However, the principal bottleneck of \WP-based inference is \emph{loop invariants}: backward propagation through a loop requires an inductive invariant that is in general not derivable from the loop body alone, and without it the rest of the inferred specification becomes vacuous.
LLMs help here: they are good at the pattern recognition loop invariants demand --- spotting that a counter is monotone, that a sum is preserved, that an element stays in a sorted prefix --- and at proposing specifications in the idiom a human would use.

Both signals are validated against the code by the prover; counter-examples feed back into the agent for refinement, and the user can intervene at any point.

\Paragraph{Contributions} While primarily a tools paper with preliminary results, this work presents several novel ideas.
As far as we know, previous work on LLM-based inference did not combine this approach with a mechanical component like \WP.
Earlier inference work is also not integrated with an interactive coding agent that lets the user contribute.
Moreover, our \WP algorithm benefits from reference elimination, as does our verification engine; this combination has not been described before, and the same approach extends to comparable verification engines for safe Rust (e.g. \cite{Aeneas}).

\Paragraph{Outline}
Sec.~\ref{sec:example} starts with a step-by-step example illustrating how the tool operates for the end user.
Sec.~\ref{sec:wp} outlines \WP analysis for Move code in the presence of Rust-style references and their elimination by the Move prover via bytecode rewriting.
In Sec.~\ref{sec:skills} we discuss the system of skills and prompts we have set up for inference and verification.
Sec.~\ref{sec:conclusion} concludes with a discussion of related and future work.


%%% Local Variables:
%%% mode: latex
%%% TeX-master: "main"
%%% End:
